\chapter{What the literature contributes}
\label{ch:literature}

The research problem crosses accounting, databases, programming languages,
econometrics and document understanding. Each field resolves part of it. The
review below distinguishes results that can be adopted directly, engineering
ideas that need adaptation, and evidence that does not establish the intended
product capability. Source versions and technical anchors are recorded in the
accompanying research ledgers. The search includes work outside the original
collection, with public-index searches and backward and forward citation
searches through 2 October 2026. It is a scoped design survey, rather than a
claim to enumerate every paper in these fields.

\section{Accounting structure and company dynamics}

\citet{nehmer1997}, especially Sections 2--5, supply a useful algebraic
starting point: balances occupy a structured state space, transactions change
that state, and aggregation relates detailed accounts to reports. This supports
a distinction between conservation constraints and permitted transactions.
An entry can conserve the accounting total and still be an inadmissible
transaction. AccDSL adopts that distinction in its posting validator. Its
particular event vocabulary, evidence model and inverse-history semantics are
new design decisions, not results established by the algebraic paper.

\citet{doshi2024} develops an accounting DSL with time, transaction logs and
correction mechanisms. The relevant material is in Chapters 4 and 7--9 of the
2024 thesis.
It demonstrates why corrections and dates belong in the language rather than
in informal comments. The thesis does not establish recovery of an entire
company history from public reports or inference over economically distinct
histories. \citet{babcinetchi2022}, pp.~342--345, offer a much smaller accounting
language and grammar proposal. It is useful as a syntax precedent, but it
provides insufficient semantics or evaluation to govern AccDSL by itself.

\citet{shahnazarian2004} connects accounting dynamics with behavioural and
econometric equations. Sections 2.3--2.8 describe the accounting relationships;
the optimisation and econometric discussion in Chapter 3 and the data and
evaluation material in Chapters 4 and 8 show how the model is used with
standardised Swedish corporate data. This is an important precedent for
simulating companies under changed conditions. Its standardised tax-data
starting point differs from heterogeneous public-report extraction, and its
working-capital funds formulation should not be imported as a general modern
cash-flow statement implementation. The transferable idea is the explicit
separation of accounting evolution and behavioural assumptions.

\citet{pierson2020}, Sections 3--4 and the reporting structures, makes that
separation operational in system dynamics. Receivables, payables and wages
accumulate differences between accrual and payment flows; reporting-period
accumulators turn instantaneous flows into statements. AccDSL can use these
structures as examples of coherent movement schedules. The model's payment
delays, inventory approximations and simplified equity treatment remain
modelling choices. They are not universal defaults for companies or evidence
that an aggregate state is sufficient for all shocks.

Bank models demonstrate why journal generation needs economic state.
\citet{halaj2016}, Section 3, combines funding shocks, fire sales, loans,
securities and risk constraints. The homogeneous portfolios and specified
funding process make the model tractable, but limit instrument-level
interpretation. \citet{behn2019}, Section 2, make managers' choices over asset
growth, debt, risk and payouts depend on capital and liquidity constraints.
These models motivate interfaces for management policy and regulatory
calculations. Their parameterisations must be calibrated and evaluated for
the target application; passing their accounting identities would not validate
an AccDSL forecast.

A stock-flow model makes the role of counterparties especially clear. The
recovered journal paper by \citet{caiani2016} describes heterogeneous loans and
capital vintages, then orders production, credit and payments within each
period (Sections 2--3). Its accounting checks precede comparisons with
empirical regularities. These distinctions support keeping economic state,
event order and accounting validation separate in AccDSL. The authors'
steady-state opening in Section 4 and Appendix A serves an artificial economy;
it cannot replace a company's observed and unresolved opening positions.
Nor can a shared label establish a shared measurement: their operating cash
flow in Section 3.1.3 subtracts loan principal repayments. The inspected
author-linked code makes that subtraction too. AccDSL must retain this
debt-service measure under a separate concept and reconcile its components
before comparing it with statement operating cash flow.

Financial planning asks a related but different question. In the retained
1974 working paper, \citet{gentry1974wp} asks what new investments must earn
to attain management's chosen earnings-per-share growth. Dividend and borrowing
decisions interact with that objective over successive years. This suggests
an explicit interface for managerial objectives and financing responses, but
does not determine customer demand or collection dates. Its independently
sampled managerial inputs are assumptions, not default probability laws. The
1979 journal version has not been checked for equivalence, and the appendix
derivation referenced in the working paper was not located in the retained
extraction. The present comparison therefore uses the planning structure,
not a reconstructed return equation.

These methods suggest a common division of work. Financial-planning equations,
stock-flow models and behavioural policies can propose economic actions.
Contract terms and dated obligations determine events; recognition and
measurement rules determine their accounting effects; reporting combines the
resulting movements. This is a provisional architectural judgment for the
initial company profile, not a finding that one simulation method dominates
the others. The Francis--Rowell simultaneous-equation alternative remains
uninspected, and a general library selection still needs exact interface and
runtime evidence.

McCarthy's 1982 REA model remains a foundational source gap for the relation
between resources, events and agents. Its uninspected technical results are
not used here. AccDSL's event and counterparty records are specified and
justified in Chapters~\ref{ch:language} and~\ref{ch:execution}. A later
multi-company equilibrium engine would require a broader stock-flow literature
and a separately justified economic model.

\section{Incomplete information and provenance}

Financial disclosures are incomplete databases in a substantive sense.
Multiple underlying records can yield the same published observations.
\citet{libkin2014}, Sections 2--6, explains the semantics of incomplete
databases, labelled nulls, conditional tables and certain answers. Repeated
occurrences of the same unknown must share a value. Conditional tables can
represent dependence between alternatives that independent missing cells
cannot express. This is directly relevant to a debt note and a cash-flow table
that refer to the same unknown repayment.

AccDSL adopts shared symbolic unknowns and conditions. It extends the database
idea with accounting constraints, continuous amounts, temporal transitions and
economic event models. A ``certain'' answer means a statement that holds for
every history in a nonempty declared feasible family. This is Boolean
certainty; a table of individually certain rows can lose existential
information and must not imply that omitted rows are absent. The payable
example in Section~\ref{sec:query-certainty} develops this distinction.
The family itself may omit
unknown mechanisms; certainty is therefore always relative to a model scope.
Imieliński and Lipski's original 1984 paper is a foundational backward lead
whose publisher copy was inaccessible in this campaign. The inspected Libkin
paper supplies the technical discussion actually used here.

\citet{green2007}, Definitions 3.1--3.2 and Section 4, treat provenance as
annotations that compose through relational operations. Combining sources
and alternative derivations can therefore have formal semantics, rather than
being a list of unrelated citations. This is useful for report aggregation and
explanation queries. Positive relational operations covered by provenance
semirings do not by themselves specify every operation AccDSL needs: negation,
nonlinear valuation, optimisation and probabilistic conditioning require their
own derivation records. Accordingly, the implementation starts with a typed
dependency graph and may use semiring annotations for supported relational
subqueries. A dependency edge is not a posterior probability.

In \emph{Temporal Databases}, \citet{snodgrass1986} distinguish historical
validity from the time facts are stored. Their snapshot, rollback, historical
and temporal examples support the need to
recover both old business facts and old database beliefs. AccDSL adds public
availability time because a database insertion date cannot establish whether
a historical investor could have seen a filing.

\section{Inverse problems, identification and probability}

\citet{stuart2010}, Section 2.2, starts from an observation equation
$y=G(u)$ and distinguishes inversion, regularisation and Bayesian inference.
The discussion makes explicit that a regularising norm or prior is an
additional assumption. AccDSL has a similar observation problem, although
its unknowns include discrete accounting interpretations and event structures
as well as continuous parameters. The infinite-dimensional well-posedness
theorems in that paper are not automatically applicable to this mixed,
potentially discontinuous model.

The relevant issue is often identification rather than computation alone.
\citet{giacomini2021}, Sections 2.1--2.6 and Theorem 1 of the archived 2020
working version, show how a class of priors can preserve uncertainty within
an identified set. A likelihood that is flat over the unresolved part does
not learn the conditional prior over that part. Posterior lower and upper
probabilities summarise a class of posterior beliefs. The paper's regularity
conditions and application to structural vector autoregressions do not prove
that an arbitrary accounting program admits the same algorithms or guarantees.
The design borrows the distinction between learnable quantities and
assumption-dependent distributions; the simple accounting bounds in
Chapter~\ref{ch:math} are derived directly.

\citet{hoeting1999}, equations (1)--(3), give Bayesian model averaging as a
weighted sum of model-conditional distributions, with weights determined by
model priors and marginal likelihoods. Sections 3 and 5 discuss computation
and prior specification. This supports an optional weighted mode, provided
AccDSL specifies the likelihood, priors and candidate model space. It does
not justify giving equal mass to every syntactically different journal
reconstruction. Two encodings of the same economic history must not gain extra
weight by duplication. Model averaging also remains conditional on the
considered model family, rather than repairing missing mechanisms.

Scenario semantics need similar care. \citet{pearl2009}, Section 3.2.1 and
equations (5)--(6), distinguishes an intervention that replaces a structural
equation from ordinary conditioning. AccDSL will distinguish a supplied
exogenous path, an observational conditional forecast and an intervention.
A rise in rates imposed by a stress test is not automatically distributed as
the collection of historical periods in which rates happened to be high.
Identifying a causal effect requires assumptions beyond reproducing accounts.

\section{Languages for rules and contracts}

\citet{vandeursen2000} survey DSL benefits, risks and implementation options.
\citet{mernik2005}, Sections 1.1 and 2.2--2.5 of the archived CWI report,
treat application libraries as serious competitors and distinguish deciding,
analysing, designing and implementing a DSL. Its implementation decision tree
explicitly includes embedding in a host language. The literature therefore
supports a conditional design decision: a DSL is attractive when domain
constraints and repeated notation justify it, but an ordinary typed library
may be cheaper and equally expressive. AccDSL will implement one shared
representation accessible from both a Python API and a textual language, so
syntax does not become a prerequisite for testing the accounting model.

\citet{merigoux2021}, Sections 3--4 and Figures 8--9, provide Catala's scopes,
exceptions, default calculus and translation into a conventional functional
target. Its explicit empty and conflict outcomes are valuable precedents for
accounting rules with exceptions. They are different from probabilistic
uncertainty. AccDSL will preserve this distinction: conflicting applicable
rules produce a diagnostic, while an evidence-supported alternative reading
produces a conditional model branch. The paper's compiler correctness results
do not establish correctness of our accounting rules or of an automatic
translation from disclosures.

The inspected LegalMath code offers reusable ideas in its shared typed
interpretation, partial observations, source references and validity/knowledge
times. Its ordinary frontend execution selects a reading. AccDSL needs an
additional execution mode that retains a family of readings, continuous
unknowns and economic state. A bounded candidate generator is useful for
search, but its candidate limit does not prove completeness.

\citet{jones2000}, Sections 3--5, build complex financial contracts from
combinators and separate their compositional meaning from a chosen valuation
model. The ACTUS technical specification similarly separates contract states,
scheduled events, state transitions and payoffs, with risk factors supplied
through an interface \citep{actus}. These are strong reasons to represent
contractual obligations independently of accounting classifications. Neither
contract cash flows nor a derivative price determine the recognition,
measurement and presentation rules of a complete IFRS or US GAAP book.

\section{Extraction and financial reasoning}

The document literature supports a staged ingestion process. PubTables-1M
\citep{smock2022}, in its canonicalisation procedure and evaluation section,
distinguishes table detection, structure recognition and functional analysis.
Its GriTS evaluation compares reconstructed table structure. TableBank
\citep{li2020}, Sections 3--4, uses document structure to obtain weak supervision
for table detection and recognition. GTE and FinTabNet \citep{zheng2021},
Sections 3--5, provide a financial-document setting and company-level dataset
splits. These sources help select extraction components and evaluation units.
Their table scores do not establish that a debt note has been interpreted
correctly across years or that every off-balance-sheet exposure has been found.

The Docling report \citep{nassar2024}, Section 3, describes PDF parsing, page
layout analysis and table structure recognition within a document-conversion
pipeline. The locally archived version lists Christoph Auer first; this
version difference is recorded in the catalogue. Docling is an ingestion
candidate whose output must retain source locations and undergo accounting
validation. Its code and model licences must be considered separately when
packaging a deployment.

\enlargethispage{\baselineskip}
TableFormer makes the division between structure and text concrete.
\citet{nassar2022tableformer}, Sections 4--5, describe a model that receives a
cropped table image and predicts its HTML structure together with a bounding
box for each cell. The cell's words and numbers come from separately supplied
PDF text; they are not generated by the structure decoder. Matching the
predicted cells to that text is itself consequential. The supplementary
algorithm can remove duplicate columns, attach previously unmatched text to a
cell, or create a new cell \citep[Section 2]{nassar2022tableformersupp}. A
plausible grid therefore still needs its text assignments and structural
corrections checked against the source.

The paper's evaluation preserves this distinction. Table 4 reports TEDS,
a similarity measure based on tree edit distance, for structure \emph{with
content} on PubTabNet PDF documents containing cropped tables: 95.4 for simple
tables, 90.1 for complex tables and 93.6 overall. The content comes from PDF
cells; Table 2 separately reports structure-only results
\citep[Tables 2 and 4, p.~4620]{nassar2022tableformer}. These are results for
that input population and text supply. They do not establish extraction
accuracy for complete financial filings, or whether a cell has been assigned
the correct entity, period, unit or accounting meaning. Those questions require
the source-linked accounting checks developed here.


FinQA \citep{chen2021}, Sections 3--6, provides questions with executable
reasoning programs and supporting evidence. TAT-QA \citep{zhu2021}, Sections
2--4, combines table cells and text spans with aggregation operations, and
explicitly treats numerical signs in evaluation. DocFinQA \citep{reddy2024}
extends financial questions to long documents and evaluates retrieval/context
choices. FinBen \citep{xie2024} broadens the collection of financial tasks and
model evaluations. These benchmarks can test retrieval and local numerical
reasoning. Their question-answer format does not require a complete
multi-period company representation, joint alternative histories or
scenario-dependent journal generation.

A directly relevant recent competitor is the auditable RAG architecture of
\citet{cosentino2026}. Sections 3.1--3.2 describe persistent extraction,
enrichment, retrieval and provenance records; Section 4.6 checks citation
resolvability and failure recovery. Section 5 limits the interpretation of
those results to the evaluated document-intelligence tasks. AccDSL should
reuse the idea of persistent evidence and recoverable processing, while adding
executable accounting and economic semantics. No performance comparison with
that architecture has been run here.

\section{Solvers and probabilistic execution}

\citet{damoura2008}, Sections 2--3, describe Z3's interfaces and theory-solving
architecture. It supports choosing an SMT solver for logical choices and
arithmetic constraints, but does not make arbitrary nonlinear simulation
tractable. The first solver fragment is deliberately bounded: finite choices,
exact linear arithmetic and selected calendar relations. Unsupported formulas
remain in the representation with a solver diagnostic.

\citet{vandeMeent2018} develop probabilistic programs through sampling,
observations and inference; the inspected material includes the first-order
language and weighted sampling discussion. Stan \citep{carpenter2017}, its
model language and sampler-output discussion, and Pyro \citep{bingham2019},
Sections 2--3, show different ways to separate a probabilistic model from
inference machinery. They are candidates for optional inference over an
already-defined economic model. Exact accounting identities and discrete rule
selection should not be turned into nearly satisfied likelihood penalties
merely to fit an inference library. Nor should successful sampler execution
be described as evidence that all interpretations were explored.

\clearpage
\section{The remaining design problem}

The combined literature provides algebra for journals, semantics for missing
information, temporal storage, provenance, contract languages, rule languages,
inference methods and extraction tools. The unresolved integration problem is
to retain their distinctions while connecting them: a disclosure constrains
economic histories; a history supplies contract and behavioural state; a
scenario drives events; accounting rules produce journals; reporting rules
render statements; and every step remains explainable.

The design is therefore an integration proposal with several local
mathematical constructions. Its originality and empirical advantage have not
been established by this survey. Its practical merit will depend on the
acceptance cases, accounting review and comparative implementation work in
Chapter~\ref{ch:validation}.
